\chapter{Hito 1: Definición del problema y determinación del alcance}

En este capítulo se detalla la definición del problema para el proyecto de la 
asignatura Desarrollo e Integración de Servicios de IA, así como la determinación de 
su alcance.

% ── 1.1 ─────────────────────────────────────────────────────────────────────
\section{Contexto y definición del problema}

En la industria vitivinícola, la certificación de la calidad y autenticidad del 
vino son procesos de gran relevancia económica, legal y reputacional. La correcta 
identificación de la variedad de uva y la detección de irregularidades en la 
composición del vino son requisitos indispensables para garantizar el cumplimiento 
de las normativas de las Denominaciones de Origen, proteger al consumidor frente a 
fraudes de etiquetado y mantener la competitividad de la bodega en un mercado 
cada vez más exigente.

Actualmente, este proceso depende casi en su totalidad de dos métodos 
complementarios: los análisis físico-químicos de laboratorio y la cata sensorial 
realizada por enólogos expertos. Aunque ambos métodos ofrecen resultados de calidad, 
presentan limitaciones evidentes cuando se aplican a gran escala. Por un lado, los 
análisis de laboratorio son precisos pero lentos y costosos. Por otro, la cata 
sensorial, aunque valiosa, introduce subjetividad y depende de la disponibilidad 
de personal especializado, lo que dificulta su aplicación sistemática en el día 
a día de una bodega con alta producción.

Además, las características del vino varían de una cosecha a otra en función del clima, 
el suelo, las técnicas de cultivo y los procesos de fermentación. Esto significa que 
los criterios de evaluación deben adaptarse continuamente, algo que los procedimientos 
manuales actuales no están diseñados para hacer de forma ágil y automatizada.

El problema que aborda el proyecto es, por un lado, la necesidad de automatizar 
y objetivar la identificación de la variedad vinícola a partir de datos medibles en 
laboratorio; por otro, la necesidad de disponer de un sistema que se mantenga 
fiable en el tiempo, capaz de detectar cuándo las características de los nuevos lotes 
de producción se alejan de los patrones conocidos y de alertar al equipo técnico para que 
tome las medidas oportunas.

En definitiva, no se trata solamente de predecir la calidad de un vino, sino de construir 
un servicio inteligente de apoyo a la decisión que acompañe al enólogo en su 
trabajo diario, aportando consistencia, transparencia y capacidad de respuesta ante los 
cambios propios del sector.

La fuente de datos seleccionada es el \textit{Wine Recognition Dataset}, 
disponible en el repositorio UCI Machine Learning Repository \cite{UCIWine}. 
Recoge análisis fisicoquímicos de vinos italianos de tres cultivares distintos, 
descritos mediante 13 propiedades químicas. La variable objetivo es la \textbf{clase de cultivo}, que toma tres valores posibles correspondientes a las tres variedades vinícolas presentes en el dataset.


% ── 1.2 ─────────────────────────────────────────────────────────────────────
\section{Objetivos generales y específicos}

El problema puede formularse con la siguiente pregunta: \textit{¿Cómo diseñar un 
sistema de inteligencia artificial capaz de identificar la variedad de un vino y 
detectar anomalías en su composición de forma objetiva, consistente y mantenible 
a lo largo del tiempo, apoyando la toma de decisiones del equipo técnico de una bodega?}

Para dar respuesta a esto, se establece como \textbf{objetivo general} del 
proyecto diseñar, implementar y desplegar un sistema de inteligencia artificial que 
permita a una bodega auditar la calidad y autenticidad de su producción vinícola de 
forma automática, objetiva y continuada, garantizando su utilidad real en un entorno 
de producción.

Los \textbf{objetivos específicos} que permitirán alcanzar este objetivo general son:

\begin{itemize}
    \item Formalizar el concepto de calidad y autenticidad vinícola desde un punto de 
    vista cuantitativo, identificando qué propiedades fisicoquímicas son relevantes 
    para la diferenciación de variedades.
    
    \item Seleccionar y procesar los datos disponibles sobre la composición química del 
    vino, asegurando su calidad y fiabilidad como base del sistema.
    
    \item Implementar y comparar distintos modelos de inteligencia artificial capaces de 
    identificar la variedad de un vino a partir de su composición, seleccionando el más 
    adecuado para su uso en producción.
    
    \item Evaluar los modelos tanto con métricas técnicas como con criterios alineados con 
    las necesidades reales del sector vitivinícola.
    
    \item Diseñar el sistema como un producto integrable en el flujo de trabajo de una bodega, 
    mantenible a lo largo del tiempo y capaz de adaptarse a los cambios de cada nueva cosecha.
    
    \item Validar el sistema con datos representativos de la producción vinícola real, 
    garantizando su utilidad práctica para el personal técnico de la bodega.
\end{itemize}

% ── 1.3 ─────────────────────────────────────────────────────────────────────
\section{Determinación del alcance}

En esta sección se determina el alcance del proyecto, considerando los componentes clave 
que definen su viabilidad y su potencial como producto real.

% ── 1.3.1 ────────────────────────────────────────────────────────────────────
\subsection{Planificación}

El proyecto se desarrollará a lo largo de cinco hitos iterativos que cubren el ciclo de vida 
completo de un proyecto de Machine Learning, desde la definición del problema hasta su 
despliegue.

% ── 1.3.2 ────────────────────────────────────────────────────────────────────
\subsection{Recursos Humanos: asignación de tareas y capacidad}

El desarrollo del sistema VinOps se llevará a cabo por un equipo de cuatro integrantes 
con roles complementarios:

\begin{itemize}
    \item \textbf{Antonio Muñoz García — Experto en dominio:} valida que los 
    resultados del sistema sean coherentes con el conocimiento 
    experto del sector vitivinícola.
    
    \item \textbf{Alejandro García Fernández — Analista de datos:} responsable de 
    la recopilación y preparación de los datos fisicoquímicos 
    del vino.
    
    \item \textbf{Pablo Parrilla Cañadas — Ingeniero de Machine Learning:} 
    diseña, entrena y evalúa los modelos de clasificación.
    
    \item \textbf{Pablo Alfonso López Fernández — Ingeniero de despliegue:} garantiza 
    que el sistema funcione de forma estable y pueda actualizarse 
    ante nuevas cosechas.
\end{itemize}

El \textbf{cliente} es una bodega o cooperativa vitivinícola, cuyo director técnico actúa 
como \textit{stakeholder} principal, validando los resultados y orientando la evolución 
del sistema.


% ── 1.3.3 ────────────────────────────────────────────────────────────────────
\subsection{Viabilidad. Estudio del ROI}

La viabilidad del proyecto se evalúa considerando tanto su impacto 
técnico como económico. Para la estimación del ROI se parten de los 
siguientes supuestos orientativos:

\begin{itemize}
    \item \textbf{Duración del proyecto:} 3 meses.
    \item \textbf{Equipo técnico:} 4 personas.
    \item \textbf{Dedicación estimada:} 10 horas semanales por persona.
    \item \textbf{Coste estimado por hora:} 35 euros.
\end{itemize}

\textit{Los valores anteriores son orientativos y podrían variar 
en función del contexto real del proyecto y del perfil del equipo 
técnico contratado.}

\[
\text{Coste total} = 4 \times 10 \; \text{h/semana} \times 
12 \; \text{semanas} \times 35 \; \text{€/hora} = 16{.}800 \; \text{€}
\]

En cuanto a los beneficios, una bodega mediana realiza aproximadamente 
2.000 análisis de calidad por campaña a un coste estimado de 15 euros 
por análisis. Automatizando el 70\% de ellos, el ahorro anual sería de 
21.000 euros. Considerando además la reutilización del sistema en 
campañas sucesivas, el ROI estimado es el siguiente:

\[
\text{ROI} = \frac{21{.}000 - 16{.}800}{16{.}800} \approx \mathbf{25\%}
\]

Este resultado indica que el proyecto es económicamente viable desde el 
primer año, con un retorno creciente en campañas posteriores.

% ── 1.3.4 ────────────────────────────────────────────────────────────────────
\subsection{Valor}

El valor aportado por el sistema VinOps se manifiesta en distintas dimensiones: 
productiva, social, organizativa y tecnológica.

\textbf{Valor productivo.} Desde el punto de vista operativo, VinOps aporta a 
la bodega un mecanismo objetivo y sistemático de clasificación de variedades vinícolas 
y detección de anomalías en la composición del producto. Este enfoque permite 
complementar el criterio del enólogo con una base cuantitativa sólida y reproducible, 
contribuyendo a identificar de forma consistente desviaciones en la calidad antes de 
que el producto llegue al mercado y a garantizar que los criterios de evaluación se 
mantengan estables a lo largo del tiempo, independientemente del personal disponible.

\textbf{Valor social y reputacional.} La certificación objetiva y transparente de la 
variedad y calidad del vino refuerza la credibilidad de la bodega ante consumidores, 
distribuidores y organismos reguladores. En un mercado donde la trazabilidad y la 
autenticidad son cada vez más valoradas, disponer de un sistema de auditoría inteligente 
puede convertirse en un argumento diferenciador frente a la competencia y en una garantía 
de cumplimiento normativo ante las Denominaciones de Origen.

\textbf{Valor organizativo.} El sistema VinOps constituye un activo estratégico reutilizable 
campaña tras campaña, extensible a nuevas líneas de producto y potencialmente integrable en 
los procesos internos de la bodega. Su diseño como producto, y no como experimento puntual, 
garantiza que la inversión realizada genere valor de forma continuada y que el sistema 
pueda evolucionar junto con las necesidades de la organización.

\textbf{Valor tecnológico.} Desde una perspectiva científico-técnica, el proyecto propone 
una metodología reproducible para la evaluación automatizada de la calidad vinícola, validada 
sobre datos reales de competición procedentes del repositorio UCI Machine Learning Repository 
\cite{UCIWine}. La combinación de modelos interpretables con criterios de evaluación 
alineados con el dominio vitivinícola sitúa a VinOps no solo como una herramienta de 
predicción, sino como una plataforma de inteligencia aplicada al sector agroalimentario, 
con capacidad de adaptación y mantenimiento a largo plazo.